% !TeX root = ../tg.tex
\chapter{文献调研}

% 近年来，随着联邦学习（Federated Learning, FL）与分割学习（Split Learning, SL）技术的不断发展，其结合形式——分割-联邦学习（SplitFed）逐渐成为分布式机器学习领域的研究热点。然而，SplitFed在理论分析、算法设计与系统优化等方面仍存在诸多挑战。以下从三个核心维度对现有文献进行系统梳理。

% \section{SplitFed框架下的收敛性分析}

% 针对传统联邦学习的收敛性分析已有较为成熟的研究成果 \cite{wang2019adaptive, li2020convergence, khaled2020tighter, karimireddy2020scaffold, koloskova2020unified, li2019convergence}。例如，Xiang Li等人在非独立同分布（Non-IID）数据条件下，证明了FL在光滑且强凸问题中的收敛时间复杂度 \cite{li2019convergence}；而Sai Praneeth Karimireddy等人的研究表明，Non-IID数据会引发客户端漂移（Client Drift），从而显著影响模型的收敛速度与稳定性 \cite{karimireddy2020scaffold}。

% 相较之下，关于SplitFed框架的收敛性分析尚处于起步阶段。据我们所知，目前仅有Li等人 \cite{li2023convergence} 对SplitFed进行了初步的理论建模与收敛性分析。该研究虽然提供了基本的收敛界，但尚未考虑实际应用中常见的标签噪声与Non-IID数据等复杂因素。因此，构建一套适用于现实场景的SplitFed收敛性分析框架，仍是亟待解决的关键问题。

% \section{应对标签噪声与Non-IID数据的SplitFed算法}

% 目前，尚未有专门针对SplitFed提出的同时具备标签噪声与Non-IID数据鲁棒性的算法。Thapa等人 \cite{thapa2022splitfed} 提出了SplitFed的基本框架，并初步验证了其在Non-IID数据上的有效性，但在处理标签噪声方面未作深入探讨。

% 在联邦学习领域，已有多种方法被用于提升模型在标签噪声或Non-IID数据下的鲁棒性。例如，Ke等人 \cite{ke2023impact} 从理论上分析了标签噪声对模型泛化能力的影响；Xu等人 \cite{b9} 提出了一种多阶段标签清洗机制；Ji等人 \cite{ji2024fedfixer} 探索了通过个性化模型辅助识别噪声标签的方法；Feddiv \cite{li2024feddiv} 引入全局噪声过滤器以提高训练稳定性；FedNoRo \cite{wu2023fednoro} 利用知识蒸馏技术识别并剔除噪声样本；Lycklama等人 \cite{lycklama2023rofl} 通过对客户端更新施加范数约束增强鲁棒性；FedLNL \cite{zhou2024federated} 结合贝叶斯推理与正则化策略提升在噪声环境下的准确性。

% 此外，针对Non-IID数据，Tian Li等人提出了FedProx \cite{li2020federated}，通过引入正则项限制本地模型偏离全局模型；Reddi等人 \cite{reddi2020adaptive} 提出了基于Adam优化器的FedAdam，利用动量信息加速收敛；Wu等人 \cite{wu2022fedadapt} 提出FedAdapt，根据客户端更新动态调整聚合策略，以适应高度异构的数据分布。

% 尽管上述方法在传统FL中取得一定成效，但由于SplitFed中模型拆分结构及客户端与服务器之间的训练耦合特性，这些算法难以直接迁移至SplitFed框架。更重要的是，它们大多仅针对单一问题（如标签噪声或Non-IID）进行优化，缺乏同时应对两类问题的能力。

% \section{低开销的训练方法：无反向传播训练}

% 当前，SplitFed普遍采用反向传播方法进行模型训练，这对边缘设备的计算资源提出了较高要求。此外，反向传播过程还需额外通信步骤来传递梯度信息，进一步增加了客户端的通信负担。

% 为降低训练开销，已有研究尝试引入无需反向传播的梯度估计方法。例如，Spall等人 \cite{spall1992multivariate} 提出的SPSA方法通过参数扰动近似梯度；Salimans等人 \cite{salimans2017evolution} 提出的进化策略（ES）利用噪声种群优化参数方向。然而，这些方法在大规模模型上表现不佳，且难以与主流深度学习框架兼容。

% 近期，Jiang等人 \cite{jiang2023one} 提出了一种新型无反向传播方法，通过注入噪声观察损失变化来估计梯度，显著降低了内存与计算需求，同时具备良好的可扩展性和框架兼容性。这一方法展现出在SplitFed中应用的潜力，有望有效缓解客户端侧的计算瓶颈和通信压力。

% 截至目前，尚未有研究将无反向传播方法应用于SplitFed框架。本课题拟首次探索此类方法在SplitFed中的可行性，旨在实现更低功耗、更轻量化、更具部署灵活性的分布式训练模式。





\section{分割-联邦学习框架下的收敛性分析}
目前，针对联邦学习的收敛性分析已有大量研究成果 \cite{wang2019adaptive, li2020convergence, khaled2020tighter, karimireddy2020scaffold, koloskova2020unified, li2019convergence}。Xiang Li等人证明了非独立同分布数据下联邦学习在光滑和强凸问题中的收敛时间复杂度\cite{li2019convergence}。
Sai Praneeth Karimireddy等人证明了非独立同分布数据会在联邦学习中引起客户端漂移（Client-Drift）从而导致
模型收敛速度慢和收敛过程不稳定\cite{karimireddy2020scaffold}。然而，关于分割-联邦学习框架下的收敛性研究却相对匮乏，据我们所知，仅有一篇文献 \cite{li2023convergence} 对该方向进行了初步探讨。现有工作尚未建立完整的收敛性分析框架，特别是在面对标签噪声和非独立同分布数据等实际场景时，缺乏系统且理论严谨的分析。
 \section{应对异构数据的分割-联邦学习算法}
 目前还没有研究提出过针对标签噪声和非独立同分布数据都具有鲁棒性的分割-联邦学习算法。目前分割-联邦学习中标签噪声的研究还较为有限，更多研究集中在非独立同分布数据上。Thapa 等人在文献 \cite{thapa2022splitfed} 中提出了分割-联邦学习算法，并探讨了其在非独立同分布数据上的鲁棒性。

联邦学习中已有很多分别应对标签噪声和非独立同分布数据的鲁棒性算法的研究。 文献\cite{ke2023impact,wu2023fednoro,b9,lycklama2023rofl, b23, ji2024fedfixer,zhou2024federated,li2024feddiv}专注于标签噪声对联邦学习的影响。Ke 等人在文献 \cite{ke2023impact} 中从理论上分析了标签噪声对联邦学习中模型泛化能力的影响。Xu 等人在文献 \cite{b9} 中提出了一种通用的联邦学习框架，通过多阶段过程检测和纠正噪声标签来应对标签噪声问题。Ji 等人在文献 \cite{ji2024fedfixer} 中探索结合全局模型与个性化模型的交替更新策略，以有效识别并过滤异构性的标签噪声。Feddiv \cite{li2024feddiv} 使用了一种全局噪声过滤器，利用来自所有客户端的信息有效识别并剔除噪声标签，从而提高训练稳定性。FedNoRo \cite{wu2023fednoro} 利用知识蒸馏技术识别噪声客户端。Lycklama 等人在文献 \cite{lycklama2023rofl} 中通过对客户端更新施加约束（如范数限制）来增强联邦学习的鲁棒性，并有效缓解标签噪声带来的影响。FedLNL \cite{zhou2024federated} 则通过交替更新本地噪声转换矩阵和分类器，并结合贝叶斯推理和本地多样性乘积正则化方法，提升了在标签噪声环境下的模型鲁棒性和准确性。
还有一些工作通过修改聚合函数构建了对非独立同分布数据具有鲁棒性的联邦学习算法。Tian Li在FedAvg的基础上引入了一个正则化项，提出了FedProx\cite{li2020federated}。
FedProx通过限制本地模型偏离全局模型的程度，帮助应对非独立同分布数据。Sashank J. Reddi等人提出了FedAdam\cite{reddi2020adaptive}，这是一种
基于Adam优化器的联邦学习算法。其使用自适应优化方法，对全局更新使用一阶和二阶动量信息，以加速收敛，对于带有非独立同分布数据的复杂任务有一定的提升效果。
Di Wu等人提出了FedAdapt\cite{wu2022fedadapt}，它会动态调整聚合策略以适应不同客户端的模型更新步幅，以更好地处理高度非独立同分布的数据。

由于分割-联邦学习中模型被拆分且客户端与服务器之间存在训练耦合，因此这些针对联邦学习设计的鲁棒性算法算法不能直接应用于分割-联邦学习的场景。同时，这些算法也无法很好地同时处理非独立同分布数据和标签噪声。
 
通常在非独立同分布数据存在的情况下，分割-联邦算法有效处理标签噪声的能力会下降，其关键原因在于：在联邦学习中，服务器无法访问客户端的数据标签，因此无法了解每个客户端的数据分布情况。而我们通过引入分割-联邦学习来解决这一问题，因为在分割-联邦学习中，客户端数据的标签存储在服务器端，使得服务器能够获取所有客户端的数据分布信息。

\section{低开销的训练方法：无反向传播训练}
目前分割-联邦学习全部使用反向传播方法训练模型，这也是目前机器学习训练模型的主流方法。但反向传播计算梯度的过程需要大量的内存和计算资源，这对于
参与分割-联邦学习的边缘设备来说是十分奢侈的。此外，分割-联邦学习中客户端与服务器需要额外进行一次通信来完成反向传播的过程，间接地增大了
客户端的通信负载。已经有很多研究提出了不需要反向传播计算梯度的方法\cite{ma2020hsic,wang2021learning,jacot2018neural,nokland2016direct,spall1992multivariate,salimans2017evolution}，其中Spall等人\cite{spall1992multivariate}提出了同时扰动随机逼近（SPSA）方法，通过在参数上施加相反符号的扰动来近似梯度。
Salimans 等人\cite{salimans2017evolution}提出了进化策略（ES）方法，通过向模型中注入随机噪声种群并沿最优方向优化参数。但这些方法都无法在较大尺寸模型上保持较高的准确率，并且这些无反向传播训练模型的方法
与通用的机器学习框架不匹配，需要额外创建新的训练流程，无法与通用的数据处理、前向传播、损失计算等步骤相匹配。文献\cite{jiang2023one}提出了一种新型的无反向传播训练模型的方法，该方法通过向模型参数或激活值中注入噪声，观察损失的变化来估计梯度，这种方法极大地
节约了设备的内存和算力，同时可以潜在地减少分割-联邦学习中客户端与服务器之间的通信负载。该方法的另一个优点是可以适配主流的机器学习框架，这
为其与分割-联邦学习相结合铺平了道路。目前，还没有分割-联邦学习算法使用无反向传播训练方法，本研究将首次使用该方案来降低客户端训练过程中的内存、计算量和通信负载。